% =====================================================================
%  CV de Jennifer Guerrero López · Versión en español
%  Compilar:  latexmk -xelatex Jennifer-Guerrero-CV-ES.tex
%  Necesita jgcv.sty y la carpeta fonts/ en la misma carpeta.
% =====================================================================
\documentclass[letterpaper]{article}
\usepackage{jgcv}

\hypersetup{
  pdftitle    = {Jennifer Guerrero López | Currículum},
  pdfauthor   = {Jennifer Guerrero López},
  pdfsubject  = {Química Farmacéutica Bióloga (QFB): formulación cosmética y soporte técnico-comercial},
  pdfkeywords = {QFB, Química Farmacéutica Bióloga, formulación cosmética, cuidado personal, cuidado del hogar, Personal \& Home Care, emulsiones, modificadores reológicos, estabilidad, viscosidad, evaluación sensorial, corneometría, PNO, I+D, soporte técnico, servicio técnico, ventas técnicas, BPL, BPD, cultivo celular, ELISA, Western blot, inmunofluorescencia, bilingüe, inglés nativo, atención a clientes, finanzas},
  pdflang     = {es-MX},
}

\cvname{Jennifer Guerrero López}
\cvheadline{Química Farmacéutica Bióloga (QFB) \cvsep Formulación cosmética y soporte técnico-comercial}
\cvcontact{Ciudad de México \cvsep \href{https://wa.me/525625903121}{+52 56 2590 3121 (solo WhatsApp)} \cvsep \href{tel:+525590027525}{+52 55 9002 7525 (casa)} \cvsep \href{mailto:jennysherly2000@gmail.com}{jennysherly2000@gmail.com}}
\cvlinks{\href{https://www.linkedin.com/in/jenn-sher}{linkedin.com/in/jenn-sher} \cvsep \href{https://jenn-sher.github.io}{Portafolio: jenn-sher.github.io}}

\begin{document}
\makecvheader

% ---------------------------------------------------------------------
\cvsection{Perfil profesional}
QFB por la UNAM. Como becaria en el laboratorio de soporte técnico de Ashland formulé productos de cuidado personal y del hogar y estudié cómo los modificadores reológicos definen la textura de una emulsión. También impartí seminarios de venta de materias primas. Inglés y español nativos, con experiencia en ventas B2B y un diplomado en finanzas de 240 horas. Busco roles que unan el laboratorio con la atención al cliente.

% ---------------------------------------------------------------------
\cvsection{Experiencia}
\cventry{Asistente becaria de I\&D}{Ene 2025\enspace–\enspace Ago 2025}
        {Ashland Chemical de México}{Personal \& Home Care Specialties (productos de cuidado personal y del hogar)}
\begin{cvitems}
  \item Formulé y reformulé productos de cuidado personal y del hogar.
  \item Estudié el efecto de dos modificadores reológicos en 14 emulsiones con un solo emulsificante.
  \item Realicé análisis cuantitativos, pruebas de estabilidad y evaluaciones sensoriales y de desempeño.
  \item Redacté el PNO de operación y limpieza del Multi Skin Test Center MC 900 (corneómetro, cutómetro y sebómetro).
  \item Impartí seminarios para impulsar la venta de materias primas y capacité a equipos de cuidado personal.
\end{cvitems}

\cventry{Asistente de laboratorio (servicio social)}{Feb 2024\enspace–\enspace Nov 2024}
        {Instituto Nacional de Rehabilitación “Luis Guillermo Ibarra Ibarra”}{Neurociencias Clínicas}
\begin{cvitems}
  \item Colaboré en un estudio de compuestos derivados de plantas en modelos in vitro de cáncer y atrofia muscular.
  \item Apliqué cultivo celular, ELISA, Western blot e inmunofluorescencia conforme a BPL y BPD.
  \item Atendí directamente a pacientes en flebotomía y en la aplicación de pruebas psicológicas.
\end{cvitems}

\cventry{Ejecutiva de ventas y atención a clientes, bilingüe}{Ene 2019\enspace–\enspace Feb 2020}
        {Teleperformance}{Cuenta de Reliant Energy (Estados Unidos)}
\begin{cvitems}
  \item Realicé ventas telefónicas B2B en inglés y en español y brindé apoyo en planta al equipo de ventas.
\end{cvitems}

% ---------------------------------------------------------------------
\cvsection{Formación}
\cventry{Licenciatura en Química Farmacéutico Biológica (QFB)}{2019–2025}
        {Universidad Nacional Autónoma de México (UNAM)}{Facultad de Química}
\cvnote{Titulación por diplomado en 2026. Título y cédula profesional en trámite.}

\cventry{Diplomado en Finanzas, 240 horas}{2026}
        {UNAM}{Facultad de Química}
\cvnote{Matemáticas financieras, amortizaciones, inversiones, nómina e ISR, ventas, mercadotecnia y administración.}

\cvline{X Simposio de Tecnología de Productos para el Cuidado Personal}{Constancia de asistencia (15 h)}{2025}
\cvline{Sociedad Estudiantil de Ciencias Químicas (SECIQ), Facultad de Química}{Miembro}{2024}

% ---------------------------------------------------------------------
\cvsection{Proyectos}
\begin{cvitems}
  \item Proyecto integrador del diplomado: creación de una empresa y gestión de sus finanzas \cvyear{2026}
  \item Gloss labial con efecto espejo: desarrollo, formulación y costeo \cvyear{2025}
  \item Texturas diferenciadas y prototipo de crema, ambos con un solo emulsificante \cvyear{2025}
  \item GlowGuard, filtro solar facial en barra sin agua, FPS 30 (proyecto en equipo) \cvyear{2024}
  \item Proyecto de emprendimiento para un laboratorio clínico \cvyear{2023}
  \item Investigación y video informativo sobre sulfonamidas \cvyear{2022}
\end{cvitems}

% ---------------------------------------------------------------------
\cvsection{Habilidades}
\cvskill{Laboratorio}{Formulación cosmética; Emulsiones y modificadores reológicos; Estabilidad; Evaluación sensorial; Análisis cuantitativo; Cultivo celular; ELISA; Western blot; Inmunofluorescencia; BPL y BPD}
\cvskill{Equipo}{Viscosímetro Brookfield; Multi Skin Test Center MC 900; Centrífuga; Estufa de estabilidad; Homogeneizador; Potenciómetro; Campana de flujo laminar}
\cvskill{Negocio}{Presentaciones técnicas; Capacitación; Ventas B2B; Atención a clientes; Finanzas; Nómina e ISR}
\cvskill{Herramientas}{Microsoft Office (Excel, PowerPoint, Word); Canva}
\cvskill{Idiomas}{Español (nativo); Inglés (nativo)}

\end{document}
